% Generated draft: requires mathematical and source review.
\documentclass[11pt]{article}
\usepackage[margin=25mm]{geometry}
\usepackage{fontspec}
\setmainfont{Latin Modern Roman}
\usepackage{amsmath,amssymb,mathtools,unicode-math}
\setmathfont{Latin Modern Math}
\usepackage{enumitem,needspace,xurl,hyperref}
\setlength{\parindent}{0pt}
\setlength{\parskip}{5pt}
\setlength{\emergencystretch}{3em}
\newcommand{\N}{\mathbb{N}}
\newcommand{\Z}{\mathbb{Z}}
\newcommand{\Q}{\mathbb{Q}}
\newcommand{\R}{\mathbb{R}}
\newcommand{\C}{\mathbb{C}}
\newcommand{\F}{\mathbb{F}}
\newcommand{\A}{\mathbb{A}}
\newcommand{\PP}{\mathbb{P}}
\newcommand{\E}{\mathbb{E}}
\newcommand{\eps}{\varepsilon}
\newcommand{\dd}{\,\mathrm{d}}
\begin{document}
\section*{UnsolvedMath ID: 30005532}
\textbf{Source problem code:} catalogue entry
\subsection{Definitions and mathematical statement}
For the purpose of explaining the surviving fragment, let \(K\) be a field and let \(W\in K[x^{\pm1},y^{\pm1}]\) be a nonzero Laurent polynomial. Write
\[
W(x,y)=\sum_{(a,b)\in S}c_{a,b}x^ay^b,
\qquad
S=\{(a,b)\in\Z^2:c_{a,b}\ne0\}.
\]
Its Newton polygon is
\[
\operatorname{Newt}(W)=\operatorname{conv}(S)\subset\R^2.
\]
This is a general definition only: the supplied record does not specify a coefficient field, coordinates, or a Laurent-polynomial model for the Floer potentials in question.

The surviving passage refers to a conjecture concerning Floer potentials and to the assertion that their Newton polygons are as predicted. The supplied curation identifies the broader subject as the Galkin--Usnich Floer-potential conjecture. However, the passage begins in the middle of a sentence and supplies neither the full conjecture nor the predicted polygons. Consequently, it does not determine a mathematical assertion with specified objects, hypotheses, and quantifiers. In particular, no equality of potentials or polygons can be stated precisely from this record alone.

The supplied record marks the statement as \texttt{unrecoverable} and its literature status as \texttt{partially\_solved}. Its dated assessment attributes to Vianna the predicted Newton-polygon distinction and infinitely many non-Hamiltonian-isotopic tori in unspecified key cases, while reporting that the full Galkin--Usnich Floer-potential conjecture is not established. These are supplied curation claims, not independently verified mathematical conclusions. The record also reports that Galkin and Mikhalkin suggested a direct tropical translation of Galkin's original construction; it does not describe that construction or translation.
\subsection{Short English statement}
The fragment concerns predicted Newton polygons of Floer potentials associated in the supplied curation with the Galkin--Usnich conjecture. It omits the relevant geometric objects, predictions, and full conjectural claim, so a precise conjecture cannot be recovered. The supplied labels \texttt{unrecoverable} and \texttt{partially\_solved} must be retained as curation metadata pending human review.
\subsection{Sources}
\begin{itemize}
\item Local registry: \texttt{data/unsolvedmath/problems.json}, with \texttt{catalog/problems.json} as the additional-record catalog; supplied permanent UnsolvedMath ID \(30005532\), exact source problem code \texttt{catalogue entry}, title ``Galkin--Usnich Floer-Potential Conjecture Fragment.'' The record lists proposed year \(2023\) and supplies no proposer.
\item Supplied report citation: \emph{Tropical Methods in Geometry} (2023), Oberwolfach Reports, pp.\ 1263--1322, DOI \texttt{10.4171/owr/2023/23}. \url{https://doi.org/10.4171/owr/2023/23}. This is also the supplied literature-source link labeled ``Oberwolfach problem report.'' The supplied verification note says that the official report identifies the fragment broadly with the Galkin--Usnich conjecture but does not state its full mathematical claim; this verification note has not been independently checked here.
\item Supplied literature source labeled ``Vianna primary paper'': \url{https://doi.org/10.1112/jtopol/jtw002}. The supplied assessment attributes the Newton-polygon distinction and infinitely many non-Hamiltonian-isotopic tori in key cases to Vianna; the record does not specify those cases or a theorem number.
\item Literature assessment date: \texttt{2026-08-22}. The supplied dated curation reports partial resolution and says that the full Galkin--Usnich Floer-potential conjecture is not established. Its background explicitly calls this a dated literature triage derived from the supplied OWR curation and requests independent verification of time-sensitive status and cited claims. This assessment is not independently verified here.
\item Attribution in the supplied fragment: Galkin's original construction and a proposed direct tropical translation by Galkin and Mikhalkin. No complete authoritative formulation of the conjecture is supplied.
\end{itemize}
\end{document}
